% !TEX program = xelatex
% !TEX root = main.tex
% Example presentation for the community-maintained CityUHK (Dongguan) theme.
\documentclass[10pt,aspectratio=169]{beamer}
\usepackage{xeCJK}
\usepackage{mathtools}
\usetheme{CityUDG}

% Portable fonts bundled with TeX Live; no macOS-only fonts are required.
\setCJKmainfont{FandolSong-Regular.otf}
\setCJKsansfont{FandolHei-Regular.otf}
\setCJKmonofont{FandolFang-Regular.otf}
\setbeamerfont{title}{size=\Large}
\setbeamerfont{subtitle}{size=\normalsize}
\title[CityUHK (Dongguan)]{Presentation Title / 演示标题}
\subtitle{Research presentation / 学术汇报}
\author[Author Name]{Author Name / 作者姓名}
\institute{City University of Hong Kong (Dongguan)\\香港城市大学（东莞）}
\date{\today}

\AtBeginSection[]{%
  \begin{frame}
    \vfill
    \begin{beamercolorbox}[sep=10pt,center]{title}
      \usebeamerfont{title}\insertsectionhead
    \end{beamercolorbox}
    \vfill
  \end{frame}
}

\begin{document}
\begin{frame}[plain]
  \titlepage
\end{frame}

\begin{frame}{Contents / 目录}
  \tableofcontents
\end{frame}

\section{Introduction / 引言}
\begin{frame}{Getting started / 开始使用}
  \begin{itemize}
    \item Edit the title, author and institute in \texttt{main.tex}.
    \item Compile with XeLaTeX; Chinese and English can be mixed directly.
    \item 修改标题、作者和正文，保留主题与图片文件即可使用。
  \end{itemize}
  \begin{block}{Design / 设计}
    Burgundy headings, orange accents and the supplied university logo.
  \end{block}
\end{frame}

\section{Examples / 示例}
\begin{frame}{Blocks / 内容块}
  \begin{block}{Main result / 主要结果}
    Present a claim, definition or key finding here.
  \end{block}
  \begin{alertblock}{Important / 重点}
    Use the orange accent to highlight important information.
  \end{alertblock}
  \begin{exampleblock}{Example / 例子}
    支持中英文混排，也可以添加自己的图表和公式。
  \end{exampleblock}
\end{frame}

\begin{frame}{Mathematics and tables / 公式与表格}
  \begin{equation*}
    \mathcal{L}(\theta)=\frac{1}{N}\sum_{i=1}^{N}
      \bigl(f_\theta(x_i)-y_i\bigr)^2
  \end{equation*}
  \vspace{0.3cm}
  \centering
  \begin{tabular}{lcc}
    \hline
    Method / 方法 & Accuracy / 准确率 & Time / 时间 \\
    \hline
    Baseline & 0.85 & 1.0 s \\
    Proposed & 0.91 & 0.8 s \\
    \hline
  \end{tabular}
  \par\medskip{\footnotesize Illustrative data only / 数据仅用于演示}
\end{frame}

\section{Conclusion / 总结}
\begin{frame}{Summary / 总结}
  \begin{enumerate}
    \item State the problem and your approach.
    \item Summarize the evidence and main result.
    \item 说明局限与下一步工作。
  \end{enumerate}
\end{frame}

\begin{frame}{Questions / 交流}
  \centering
  {\color{DGBurgundy}\Huge Thank you! / 谢谢！}
\end{frame}
\end{document}
